\section{Limitations}
\label{sec:limits}

Our study has these limits:
\begin{itemize}
  \item \textbf{Small test sets.} With 40--80 tasks, few pairs of variants are clear. Thus, each
  count of decisions has a large error.
  \item \textbf{Mutants give most clear decisions.} A mutant is in 17 of the 20 clear pairs of W1,
  all 13 of W2 and 6 of 8 of the W3 planner. Thus, our results show that miniatures
  find large errors. Small prompt changes need larger test sets.
  \item \textbf{One small model, an LLM judge.} All workflows use one small model, and only W1 also
  ran on a second family. Rules and a model, not people, gave the grades and the reasons.
  \item \textbf{Free tools.} Our tools are free and fast, so simulated tools cost
  1.05--1.23$\times$ the full run. With paid or slow tools, simulation saves money. Its validity
  does not depend on the price of the tool.
  \item \textbf{Open tests.} We did not test a grounded simulator on W2, and W2 has no noise
  ceiling.
  \item \textbf{Overfit.} A team that tunes against one miniature for a long time can learn its
  errors. Checks on the full workflow stay necessary.
\end{itemize}
